% !TeX root = ../../main.tex
\subsection{助動詞とは}
    助動詞は動詞を助ける働きがあり、原則\fitblankbf[]{動詞の直前}に置く。

    助動詞の後ろは必ず\fitblankbf[]{動詞の原形}にする。

\subsection{疑問文と否定文}
    助動詞を含む文の疑問文は文頭に\fitblankbf[]{助動詞}を出し、否定文は助動詞の直後に\fitblankbf[]{not}をつける。
    つまり、\fitblankbf[]{be 動詞}の用法と同じである。

    \begin{theorembox}{\textbf{助動詞の用法}}
        \begin{enumerate}
            \item 助動詞は原則\fitblankbf[]{動詞の直前}に置く
            \item 助動詞の後ろは必ず\fitblankbf[]{動詞の原形}にする
            \item 助動詞を含む文の疑問文は、\fitblankbf[]{助動詞}を\fitblankbf[]{文頭}に出す
            \item 助動詞を含む文の否定文は、\fitblankbf[]{助動詞}の直後に\fitblankbf[]{not}をつける

            したがって、助動詞を含む文の用法は \fitblankbf[]{be動詞}を含む文の用法と同じである。
        \end{enumerate}
    \end{theorembox}
    \vspace{1em}

    % \begin{itemize}
    %     \item He can play \underline{the guitar} very well.
    %     \par\vspace{0.3em}
    %     \noindent\makebox[12em][l]{疑問文へ}\ $\longrightarrow$\ \dialogueblankinline{Can he play the guitar very well?}
    %     \par\vspace{0.5em}
    %     \noindent\makebox[12em][l]{否定文へ}\ $\longrightarrow$\ \dialogueblankinline{He can't play the guitar very well.}
    %     \par\vspace{0.5em}
    %     \noindent\makebox[12em][l]{下線部を問う疑問文へ}\ $\longrightarrow$\ \dialogueblankinline{What can he play very well?}
    %     \par\vspace{0.5em}
    %     \item She will visit \underline{Kyoto} next week.
    %     \par\vspace{0.3em}
    %     \noindent\makebox[12em][l]{疑問文へ}\ $\longrightarrow$\ \dialogueblankinline{Will she visit Kyoto next week?}
    %     \par\vspace{0.5em}
    %     \noindent\makebox[12em][l]{否定文へ}\ $\longrightarrow$\ \dialogueblankinline{She won't visit Kyoto next week.}
    % \end{itemize}

\newpage
